\scanpage{181}
\mspage{\deleted{\unsure{129}}\unsure{130}}
Si en particulier on a un diagramme \uline{commutatif} dans $\mathit{Cat}$
\diagram{\begin{tikzcd}[ampersand replacement=\&,row sep=small,column sep=small]
\&\& X_0 \arrow[dl,"\alpha_1"']\arrow[dr,"\alpha_2"]\&\\
Q={} \& X_1\arrow[dr,-,"f_1"']\&\& X_2\arrow[dl,"f_2"]\\
\&\&Y\&{,}
\end{tikzcd}}
on peut le compléter par \deleted{\unsure{$f_0$}}.
\[
f_0\overset{\mathrm{def}}{=}f_1\alpha_1=f_2\alpha_2
\]
en un diagramme
\diagram{\begin{tikzcd}[ampersand replacement=\&,row sep=small]
\&X_0\arrow[dl,"\alpha_1"']\arrow[dr,"\alpha_2"]\arrow[dd,"f_0", ""{name=C,near end}]\&\\
X_1\arrow[dr,"f_1"',""{name=L,near end}]\&\& X_2\arrow[dl,-,"f_2", ""{name=R,near end}]\\
\&Y\&{,}\\
\arrow[from=L,to=C,bend left=12,"\mathrm{id}"]
\arrow[from=R,to=C,bend right=12,"\mathrm{id}"']
\end{tikzcd}}
ce qui montre qu'on a une donnée \inserted{du type \unsure{$(\ast)$|cas|carré} ci-dessus,} définissant \inserted{donc} une flèche de
\diagram{\begin{tikzcd}[ampersand replacement=\&,row sep=tiny,column sep=small]
\& X_1\&\\
\unsure{\varphi|\rho}_Q:\quad\mathscr X=\int X_0\arrow[ru]\arrow[rd]\&{}\arrow[r]\&Y\\
\& X_2\&
\end{tikzcd}}
qui donne lieu à un diagramme cubique
% archive sheet number: 181
